% ~700 words. Task taxonomy, five-task descriptions, leaderboard
% mechanics, PR-extension protocol.
\section{Benchmarks}
\label{sec:benchmarks}

\subsection*{Task taxonomy}

Every FLAIR task is declared on three axes in a single \texttt{META}
dictionary: \emph{task class} (\textit{outcome}, \textit{deterioration},
\textit{intervention}, or \textit{operational}); \emph{temporal
pattern} (\textit{episode}, \textit{continuous}, or \textit{fixed
horizon}); and \emph{prediction unit} (\textit{one-per-stay} or
\textit{one-per-window}). The
triple determines the shape of the cohort and prediction parquets,
which fields are mandatory, and which evaluation primitives the report
builder invokes. New task authors fill the triple before writing any
SQL or feature code; reviewers consume it before opening a single line
of the implementation. The taxonomy is deliberately small so that the
boundary between ``new task'' and ``existing task'' is unambiguous.

\subsection*{Initial five-task suite}

The v1 release covers three outcome targets and two
intervention-anchored targets spanning ICU operations, disposition,
and infectious physiology.

\begin{itemize}[leftmargin=1.5em]
  \item \textbf{T1, ICU Daily In-Hospital Mortality (07:00 local).}
        One-per-window, 24-hour input window anchored at daily 07:00.
        Cohort: first ICU stay with discharge in
        \{long-term acute care, mortality\}. Label: in-hospital death
        or transfer to hospice during that ICU stay.
  \item \textbf{T2, ICU Daily Long-Term Acute Care Discharge
        (07:00 local).} Complementary label to T1 on the same daily
        grid: prediction that the patient will be discharged to
        long-term acute care rather than dying or being discharged
        home. Together T1 and T2 capture the disposition decision
        that anchors much downstream ICU planning.
  \item \textbf{T3, Extubation Failure within 24~h.} One-per-stay,
        anchored at the first confirmed invasive-mechanical-ventilation
        extubation event identified by the consensus-window episode
        detector. Label: death or reintubation within 24~hours.
        Benchmark targets the classic respiratory liberation question
        \citep{thille2011extubation}.
  \item \textbf{T4, Sepsis (antibiotic-anchored) within 6~h.}
        One-per-window on an hourly grid; $t_{\mathrm{zero}}$ is the
        first qualifying systemic intravenous antibiotic. Cohort follows
        the Sepsis-3 hospital-wide pattern \citep{singer2016sepsis3,seymour2016assessment,rhee2019sepsis};
        leakage prevention drops decision points at or after
        $t_{\mathrm{zero}}$. The implementation mirrors the
        antibiotic-anchored evaluation framing of
        \citet{kamran2024sepsis}.
  \item \textbf{T5, Unplanned ICU Readmission.} One-per-stay,
        anchored at first-ICU discharge; the label is a return to the
        ICU later in the same hospitalization through a non-procedural
        path. Working from the CLIF \texttt{ADT} location stream,
        planned ICU $\to$ procedural $\to$ ICU excursions are merged and
        excluded; the cohort drops death-in-ICU, direct ICU discharge,
        sub-hour first stays, and intermediate ER visits. The framing
        follows structured-EHR ICU-readmission prediction
        \citep{rojas2018readmission}.
\end{itemize}

Each cohort parquet carries a fixed ten-column schema, identifiers,
window endpoints, label, split, age, sex, race, and ethnicity, so
that downstream feature engineering and evaluation code is shared
across tasks.

\subsection*{Count-based reference baselines}

The leaderboard is anchored by a count-featurization baseline in the
EHRSHOT/FEMR tradition \citep{wornow2023ehrshot,steinberg2021femr},
adapted to CLIF. Every row of a CLIF event table (\texttt{vitals},
\texttt{labs}, \texttt{medication\_admin\_*},
\texttt{respiratory\_support}) is mapped to a deterministic
\texttt{<table>.<category>} token (e.g., \texttt{vitals.heart\_rate},
\texttt{labs.creatinine}), and each cohort window is reduced to a
sparse integer count vector over those tokens. The featurizer is
config-free and exhaustive, every schema-locked category flows
through with no per-task curation, so the same code runs unchanged on
MIMIC-CLIF and at every consortium site. Two simple heads consume the
count vectors: a single XGBoost head, emitting a probability score
\texttt{y\_prob}. It is trained on MIMIC-CLIF
and shipped frozen to sites; hyperparameter grids, the token
dictionary, and the missingness convention (absent token $=$ zero
count) are given in the supplementary appendix
(Sec.~\ref{sec:supp}). Internal-test performance on MIMIC-CLIF is
summarized in Table~\ref{tab:baseline}.

\begin{table}[t]
\centering
\caption{XGBoost count-featurization baseline on the MIMIC-CLIF
internal test split (75/25 hold-out on stitched encounters,
\texttt{hospitalization\_join\_id}, so an encounter never spans train and
test). AUROC shown with 1{,}000-iteration bootstrap 95\% CI.}
\label{tab:baseline}
\begin{tabular}{l c c}
\toprule
Task & AUROC (95\% CI) & AUPRC \\
\midrule
T1, ICU daily mortality       & 0.790 (0.784--0.795) & 0.502 \\
T2, LTAC discharge            & 0.757 (0.751--0.762) & 0.377 \\
T3, Extubation failure 24\,h  & 0.701 (0.672--0.729) & 0.124 \\
T4, Sepsis 6\,h               & 0.812 (0.792--0.854) & 0.021 \\
T5, Unplanned ICU readmission & 0.645 (0.628--0.660) & 0.252 \\
\bottomrule
\end{tabular}
\end{table}

\subsection*{Leaderboard mechanics}

For each task and site, FLAIR reports area under the receiver-operating
characteristic curve (AUROC), area under the precision--recall curve
(AUPRC), and F$_1$ at the operating threshold, each with 1{,}000-iteration
bootstrap 95\% confidence intervals. Calibration is summarized by
slope, intercept, and Brier score, with raw reliability points
downsampled into \texttt{curves.json} so that the central dashboard can
render reliability and ROC/PR curves without re-touching patient
data. Fairness is reported as per-subgroup AUROC across sex, race,
ethnicity, and age bins together with explicit parity gaps. The
leaderboard ranks submissions on the primary metric for each task and
displays per-site variance prominently: a single high pooled AUROC
hides exactly the failure mode FLAIR exists to surface.

\subsection*{An open, extensible task catalog}

FLAIR is released open source under the Apache~2.0 license at
\texttt{github.com/clif-consortium/FLAIR}, and the five tasks above
are explicitly the \emph{seed} of a growing catalog rather than the
benchmark's permanent scope. Any researcher or clinician can propose a
new task as a single Python module under
\texttt{flair\_benchmark/tasks/} containing the \texttt{META}
dictionary and a \texttt{build} function returning the canonical
cohort frame; no registry edits or library changes are required.
Acceptance review verifies (i) reproducibility against MIMIC-CLIF,
(ii) clinical face validity through a clinician-led review, and
(iii) site-governance acceptance at a minimum of two consortium sites
before the task enters the federated suite. This is the structural
counter to benchmark saturation: the catalog grows with community
need, distributed across many proposers, rather than ossifying around
the priors of any single author group
\citep{recht2019imagenet,chekroud2024illusory}.
